\chapter{\parbox{0.75\textwidth}{Оператор циклу \texttt{for} та функція \texttt{range()}}}
	
Оператори циклів використовують у випадках, коли в програмі необхідно багаторазово виконати певну послідовність команд. У мові Python основним циклом із заздалегідь відомою кількістю повторень є цикл \texttt{for}.

\vspace{0.3cm}

\begin{pydef}
Синтаксис циклу for
for параметр in range(start, stop, step):
    блок_операторів
\end{pydef}

\vspace{0.3cm}

\section*{1. Генератор послідовностей \texttt{range()}}

Найчастіше для формування послідовностей чисел у циклі \texttt{for} використовується вбудована функція \texttt{range()}. Вона може приймати від 1 до 3 аргументів:

\begin{itemize}
	\item \texttt{range(stop)} "--- генерує числа від $0$ до $stop - 1$ із кроком $1$.
	\item \texttt{range(start, stop)} "--- генерує числа від $start$ до $stop - 1$ із кроком $1$.
	\item \texttt{range(start, stop, step)} "--- генерує числа від $start$ до $stop - 1$ із кроком $step$.
\end{itemize}

\begin{table}[h!]
	\centering
	\renewcommand{\arraystretch}{1.3}
	\begin{tabular}{|l|l|l|}
		\hline
		\textbf{Виклик \texttt{range()}} & \textbf{Згенерована послідовність} & \textbf{Примітка} \\ \hline
		\texttt{range(5)} & \texttt{0, 1, 2, 3, 4} & Від 0 до 4 \\ \hline
		\texttt{range(1, 6)} & \texttt{1, 2, 3, 4, 5} & Від 1 до 5 \\ \hline
		\texttt{range(2, 10, 2)} & \texttt{2, 4, 6, 8} & Парні числа від 2 до 8 \\ \hline
		\texttt{range(10, 0, -1)} & \texttt{10, 9, 8, 7, 6, 5, 4, 3, 2, 1} & Зворотний відлік \\ \hline
	\end{tabular}
	\caption{Приклади роботи функції \texttt{range()}}
\end{table}

\begin{py}{Якщо передати в range() лише одне число, за замовчуванням відлік починається з 0, а верхня межа не включається}
for i in range(5):
    print(i)
\end{py}
Результат виведення : 0, 1, 2, 3, 4


\begin{py}{range(start, stop) починає з числа start (включно) і зупиняється перед stop (не включаючи його).}
for i in range(1,5):
    print(i)
\end{py}
Результат виведення: 1, 2, 3, 4

\begin{py}{Третій параметр — це крок (step). range(1, 10, 2) починає з 1 і додає 2 на кожному кроці, генеруючи непарні числа до 10}
for i in range(1,10,2):
    print(i)
\end{py}
Результат виведення: 1, 3, 5, 7, 9

\begin{py}{Якщо крок від'ємний (-1), значення будуть зменшуватися. Цикл починається з 10 і зупиняється перед 0.}	
for i in range(10,0,-1):
    print(i)
\end{py}
Результат виведення: 10, 9, 8, 7, 6, 5, 4, 3, 2, 1

\section*{2. Базові алгоритми з циклом \texttt{for}}

\begin{py}{Приклад 1. Обчислення суми натуральних чисел від 1 до n}
n = int(input("Введіть n "))
s = 0
for i in range(1, n + 1):
    s += i

print(f"Сума чисел від 1 до {n} дорівнює: {s}")
\end{py}

input(...) запитує у користувача число та приймає його як рядок (текст).

int(...) перетворює введений текст у ціле число та зберігає його в змінну $n$.

$s = 0$ cтворюється змінна-накопичувач $s$ із початковим значенням 0. У ній буде зберігатися підсумкова сума.

for i in range(1, n + 1): цикл for перебирає числа по порядку.

range(1, n + 1) генерує послідовність від 1 до $n$ включно (верхня межа $n + 1$ не включається в діапазон, тому цикл зупиняється саме на n).

s += i  на кожній ітерації циклу поточне число i додається до суми s (це скорочений запис для $s = s + i).$

print(f"Сума чисел від 1 до {n} дорівнює: {s}")

За допомогою $f$-рядка на екран виводиться підсумковий результат.
\begin{py}{Приклад 2. Обчислення факторіала n!}
n = int(input("Введіть n "))
fact = 1

for i in range(1, n + 1):
    fact *= i

print(f"{n}! = {fact}")
\end{py}

\begin{py}{Приклад 3. Табулювання функції y = sin(x) на відрізку [a, b]}
import math

a = float(input("Введіть початок відрізка a "))
b = float(input("Введіть кінець відрізка b "))
n = 10
h = (b - a) / n

print(" x | sin(x)")
print("-------------")
for i in range(n + 1):
    x = a + i * h
    y = math.sin(x)
    print(f"{x:6.2f} | {y:6.4f}")
\end{py}


Форматоване виведення (виділений рядок):

$\{x:6.2f\}$ — виводить значення $x$: занімає 6 символів у ширину та округлюється до 2 знаків після коми.| — виводить розділювальну риску між стовпчиками.

$\{y:6.4f\}$ — виводить значення $y$ ($\sin(x)$): шириною в 6 символів та округлене до 4 знаків після коми.

\vspace{0.5cm}

\section*{3. Аналіз дільників числа та перевірка на простоту}

\begin{tcolorbox}[colback=pyblue!5!white,colframe=pyblue,title=\textbf{Математична довідка},fonttitle=\sffamily\bfseries]
	\textbf{Просте число} "--- це натуральне число, більше за 1, яке має точно два різні натуральні дільники: 1 та саме це число. Якщо число $n$ має дільник, відмінний від 1 та $n$, то найменший такий дільник не перевищує $\lfloor\sqrt{n}\rfloor$.
\end{tcolorbox}

\begin{py}{Приклад 4. Виведення всіх дільників числа}
n = int(input("Введіть натуральне число n "))
count=0
for i in range(1, n + 1):
    if n % i == 0:
        print(i)
    count=count+1
print(f"Кількість дільників: {count}")
\end{py}

Пошук дільників: Програма перебирає всі цілі числа $i$ від $1$ до $n$ включно і перевіряє умову n % i == 0 (чи ділиться $n$ на $i$ без залишку).

Виведення та лічильник: Якщо число $i$ є дільником:воно одразу виводиться на екран (print(i)); лічильник count збільшується на $1$.  Наприкінці програма виводить загальну кількість знайдених дільників.



\section*{4. Обчислення нескінченних математичних рядів}

\subsection*{4.1. Сума геометричного ряду}
Обчислимо суму $S = \sum_{n=0}^{N} x^n = 1 + x + x^2 + \dots + x^N$ ($|x| < 1$).

\begin{py}{Приклад 5. Сума геометричного ряду}
x = float(input("Введіть x (|x| < 1) "))
n = int(input("Введіть кількість членів N "))

s = 0.0
term = 1.0

for i in range(n):
    s += term
    term *= x

print(f"Сума ряду: {s:.6f}")
print(f"Точна сума 1/(1-x): {1 / (1 - x):.6f}")
\end{py}

\subsection*{4.2. Обчислення експоненти $e^x$ через ряд Тейлора}

\begin{py}{Приклад 6. Обчислення e^x через ряд Тейлора}
import math

x = float(input("Введіть x "))
n = int(input("Кількість членів ряду "))

s = 0.0
term = 1.0

for i in range(1, n + 1):
    s += term
    term *= x / i

print(f"Наближене значення e^{x} = {s:.8f}")
print(f"Точне значення (math.exp): {math.exp(x):.8f}")
\end{py}
$f$"..." ($f$-рядок) — дозволяє вставляти значення змінних безпосередньо всередину тексту за допомогою фігурних дужок {}.

$\{s:.8f\}$ — бере значення змінної $s$ (суму ряду) та форматує його як число з плаваючою крапкою, залишаючи рівно 8 знаків після коми.

Цей рядок показує наближене значення експоненти $e^x$ , вираховане через суму перших $n$ членів ряду Тейлора.
\section*{5. Управління циклом: \texttt{break}, \texttt{continue} та \texttt{else}}

\begin{itemize}
	\item \texttt{break} "--- достроково припиняє виконання циклу.
	\item \texttt{continue} "--- пропускає залишок тіла циклу та переходить до наступної ітерації.
	\item \texttt{else} "--- блок коду, який виконується після завершення циклу, якщо вихід з нього \textbf{не} відбувся за допомогою \texttt{break}.
\end{itemize}

\begin{py}{Приклад 7. Пошук першого від'ємного числа з використанням break та else}
for i in range(5):
    val = int(input("Введіть число "))
    if val < 0:
        print(f"Знайдено перше від'ємне число: {val}")
        break
else:
    print("Усі введені числа є невід'ємними.")
\end{py}

\vspace{0.3cm}

\section*{6. Оператори інкременту та особливості \texttt{float}}

У Python відсутні оператори \texttt{++} та \texttt{--}. Замість них використовуються комбіновані оператори присвоєння:
\begin{itemize}
	\item \texttt{x += 1} замість \texttt{x = x + 1}
	\item \texttt{x -= 1} замість \texttt{x = x - 1}
	\item \texttt{x *= 2} замість \texttt{x = x * 2}
\end{itemize}

\begin{tcolorbox}[colback=yellow!10!white,colframe=orange,title=\textbf{Обережно з дійсними числами (float)!},fonttitle=\sffamily\bfseries]
	Числа з плаваючою крапкою зберігаються в пам'яті комп'ютера у двійковій системі з обмеженою точністю (53 біти). Через це обчислення можуть містити незначну похибку:
\begin{py}{Особливості float}
>>> 0.1 + 0.1 + 0.1
0.30000000000000004
\end{py}
	Для порівняння дійсних чисел завжди використовуйте функцію \texttt{math.isclose(a, b)} або порівняння з точністю $\varepsilon$: \texttt{abs(a - b) < 1e-9}.
\end{tcolorbox}

\vspace{0.3cm}

\
\end{enumerate}
